\documentclass[aspectratio=169, xcolor={table,xcdraw}]{beamer}

\usepackage[T1]{fontenc}
\usepackage[utf8]{inputenc}
\usepackage{tgpagella}
\usepackage{tgcursor}
\usefonttheme{serif}
\usepackage[spanish, es-noshorthands]{babel}
\renewcommand{\tablename}{Tabla}

\usepackage{amsmath, amssymb, bm}
\usepackage{graphicx}
\usepackage{booktabs}
\usepackage[most]{tcolorbox}
\usepackage{tikz}
\usepackage{pgfplots}
\pgfplotsset{compat=1.18}
\usetikzlibrary{shapes.geometric, arrows.meta, positioning, calc, angles, quotes}

\usetheme{Madrid}
\usecolortheme{default}
\definecolor{ucbblue}{RGB}{0, 51, 102}
\setbeamercolor{structure}{fg=ucbblue}
\setbeamercolor{title}{fg=white, bg=ucbblue}
\setbeamercolor{frametitle}{fg=white, bg=ucbblue}

\title[Dinámica Esencial]{Dinámica Esencial de Manipuladores e Integración Hito 2}
\subtitle{De la Trayectoria al Esfuerzo de los Actuadores}
\author[B. Quiroga Turdera]{M.Sc. Ing. Bernardo Quiroga Turdera}
\institute[UCB Tarija]{Universidad Católica Boliviana ``San Pablo'' — Sede Tarija}
\date{Semana 9 - Sesión 2}

\begin{document}

\begin{frame}
    \titlepage
\end{frame}

\begin{frame}{Cinemática vs. Dinámica[cite: 14]}
    ¿Por qué no es suficiente la cinemática para controlar un robot físico?[cite: 14]

    \begin{columns}[T]
        \begin{column}{0.48\textwidth}
            \begin{tcolorbox}[colback=ucbblue!5, colframe=ucbblue, title=Cinemática (LSPB), left=1mm, right=1mm, top=1mm, bottom=1mm]
                \small
                Proporciona las cantidades geométricas del movimiento[cite: 14]:
                \begin{center}
                    $q(t), \quad \dot{q}(t), \quad \ddot{q}(t)$
                \end{center}
                Describe \textbf{cómo} debe moverse el robot en el tiempo, pero ignora las fuerzas necesarias.
            \end{tcolorbox}
        \end{column}
        \begin{column}{0.48\textwidth}
            \begin{tcolorbox}[colback=red!5, colframe=red!70!black, title=Dinámica, left=1mm, right=1mm, top=1mm, bottom=1mm]
                \small
                Relaciona el movimiento con el esfuerzo de los actuadores[cite: 14]:
                \begin{center}
                    $\boxed{(q, \dot{q}, \ddot{q}) \longrightarrow \tau}$
                \end{center}
                Determina el \textbf{torque requerido} considerando masas, inercias y gravedad.
            \end{tcolorbox}
        \end{column}
    \end{columns}
    
    \vfill
    \begin{alertblock}{Ejemplo de la limitación cinemática}
        \small Dos movimientos pueden seguir la misma ruta geométrica desde $q_A$ a $q_B$. Sin embargo, si uno es muy rápido, requiere una mayor aceleración $\ddot{q}$, exigiendo un torque inercial muy diferente[cite: 14].
    \end{alertblock}
\end{frame}

\begin{frame}{Fundamentos de Euler-Lagrange[cite: 14]}
    La energía proporciona una ruta sistemática desde el modelo mecánico hasta las ecuaciones de movimiento del manipulador[cite: 14].
    
    \begin{columns}[T]
        \begin{column}{0.5\textwidth}
            \textbf{Energías del Sistema:}
            \begin{itemize}
                \item $T(q, \dot{q})$: Energía Cinética[cite: 14].
                \item $V(q)$: Energía Potencial[cite: 14].
            \end{itemize}
            
            \vspace{0.3cm}
            \textbf{Lagrangiano:}
            \begin{equation*}
                \boxed{L(q, \dot{q}) = T(q, \dot{q}) - V(q)} \text{[cite: 14]}
            \end{equation*}
        \end{column}
        \begin{column}{0.5\textwidth}
            \textbf{Ecuación de Euler-Lagrange:}
            Para cada coordenada generalizada $q_i$[cite: 14]:
            \begin{center}
                $\displaystyle \boxed{ \frac{d}{dt}\left(\frac{\partial L}{\partial \dot{q}_i}\right) - \frac{\partial L}{\partial q_i} = \tau_i }$ \text{[cite: 14]}
            \end{center}
            
            \vspace{0.2cm}
            \footnotesize $\tau_i$ representa el \textbf{torque generalizado} asociado al grado de libertad $i$ (es decir, el esfuerzo del motor en la junta)[cite: 14].
        \end{column}
    \end{columns}
\end{frame}

\begin{frame}{Derivación de Dinámica 1R (Péndulo Rígido) (1/2)[cite: 14]}
    Consideremos un eslabón rígido girando en un plano vertical[cite: 14].
    \begin{columns}[T]
        \begin{column}{0.6\textwidth}
            \textbf{Parámetros del modelo[cite: 14]:}
            \begin{itemize}
                \item $q$: ángulo medido desde la vertical hacia abajo.
                \item $m$: masa total del eslabón.
                \item $\ell_c$: distancia desde la junta al centro de masa (COM).
                \item $I_c$: inercia del COM respecto al eje de rotación.
            \end{itemize}
            
            La inercia equivalente de rotación respecto a la junta (Teorema de Steiner) es[cite: 14]:
            \begin{equation*}
                J = I_c + m\ell_c^2 \text{[cite: 14]}
            \end{equation*}
        \end{column}
        \begin{column}{0.4\textwidth}
            \centering
            \begin{tikzpicture}[scale=1.2, thick]
                \draw[dashed, gray] (0,0) -- (0,-2.5) node[below]{Vertical};
                \draw[ucbblue, line width=4pt] (0,0) -- (1.5,-2);
                \filldraw[black] (0,0) circle (2pt) node[above right]{Junta ($q=0$)};
                \filldraw[red] (0.75,-1) circle (2pt) node[right]{COM ($m$)};
                
                \draw[-Stealth, black] (0,-1) arc (270:307:1) node[midway, below]{$q$};
                \draw[-Stealth, red, line width=1.5pt] (0.75,-1) -- (0.75,-2) node[right]{$mg$};
                
                \draw[dashed] (0,0) -- (0.75,-1) node[midway, above left]{$\ell_c$};
                \draw[-Stealth, orange, line width=1.5pt] (-0.3, 0) arc (180:90:0.3) node[above left]{$\tau$};
            \end{tikzpicture}
        \end{column}
    \end{columns}
\end{frame}

\begin{frame}{Derivación de Dinámica 1R (Péndulo Rígido) (2/2)[cite: 14]}
    \begin{columns}[T]
        \begin{column}{0.48\textwidth}
            \textbf{1. Energías[cite: 14]:}
            
            $T = \frac{1}{2}J\dot{q}^2$
            
            $V = mg\ell_c(1 - \cos q)$
            
            \vspace{0.3cm}
            \textbf{2. Lagrangiano[cite: 14]:}
            
            $L = \frac{1}{2}J\dot{q}^2 - mg\ell_c(1 - \cos q)$
            
            \vspace{0.3cm}
            \textbf{3. Derivadas Parciales[cite: 14]:}
            
            $\frac{\partial L}{\partial \dot{q}} = J\dot{q} \implies \frac{d}{dt}\left(\frac{\partial L}{\partial \dot{q}}\right) = J\ddot{q}$
            
            $\frac{\partial L}{\partial q} = -mg\ell_c\sin q$
        \end{column}
        \begin{column}{0.5\textwidth}
            \textbf{4. Resultado Euler-Lagrange[cite: 14]:}
            
            Aplicando la ecuación general:
            \begin{equation*}
                J\ddot{q} - (-mg\ell_c\sin q) = \tau
            \end{equation*}
            \begin{equation*}
                \boxed{\tau = J\ddot{q} + mg\ell_c\sin q} \text{[cite: 14]}
            \end{equation*}
            
            \vspace{0.2cm}
            \begin{tcolorbox}[colback=yellow!10, colframe=orange, title=Comprobación Estática]
                \small Si el brazo no se mueve ($\dot{q}=0, \ddot{q}=0$):
                
                $\tau = mg\ell_c\sin q$
                
                En $q=0 \implies \tau=0$.
                
                En $q=\pi/2 \implies \tau=mg\ell_c$[cite: 14].
            \end{tcolorbox}
        \end{column}
    \end{columns}
\end{frame}

\begin{frame}{Ecuación Dinámica General del Manipulador[cite: 14]}
    \textbf{Requisito de Diseño para Hito 2:} Para un robot de múltiples GDL, la ecuación se generaliza matricialmente[cite: 14]:
    \begin{center}
        $\displaystyle \boxed{ \tau = M(\mathbf{q})\ddot{\mathbf{q}} + C(\mathbf{q}, \dot{\mathbf{q}})\dot{\mathbf{q}} + \mathbf{g}(\mathbf{q}) }$ \text{[cite: 14]}
    \end{center}
    
    Para un robot de $n$ grados de libertad[cite: 14]: $\mathbf{q}, \dot{\mathbf{q}}, \ddot{\mathbf{q}}, \tau, \mathbf{g} \in \mathbb{R}^n$, y $M, C \in \mathbb{R}^{n \times n}$.
    
    \vspace{0.2cm}
    \begin{center}
        \renewcommand{\arraystretch}{1.3}
        \footnotesize
        \begin{tabular}{|l|l|p{6cm}|}
        \hline
        \rowcolor{gray!20}\textbf{Término} & \textbf{Rol} & \textbf{Significado Físico} \\ \hline
        $M(\mathbf{q})\ddot{\mathbf{q}}$ & Inercial & Torque requerido para acelerar las masas del mecanismo[cite: 14]. \\ \hline
        $C(\mathbf{q}, \dot{\mathbf{q}})\dot{\mathbf{q}}$ & Velocidad & Efectos centrífugos y de Coriolis (acoplamiento entre juntas)[cite: 14]. \\ \hline
        $\mathbf{g}(\mathbf{q})$ & Gravedad & Torque requerido para sostener el robot contra la gravedad[cite: 14]. \\ \hline
        $\tau$ & Esfuerzo & Vector de torque total requerido en los actuadores[cite: 14]. \\ \hline
        \end{tabular}
    \end{center}
    \tiny{\textit{Nota: Nuestro modelo simple 1R no presentaba la matriz $C$ de Coriolis, pero está presente en casi cualquier movimiento multi-articular espacial[cite: 14].}}
\end{frame}

\begin{frame}{Integración: LSPB hacia Torque (1/2)[cite: 14]}
    La trayectoria LSPB nos suministra perfiles temporales precisos[cite: 14]: $q(t), \dot{q}(t), \ddot{q}(t)$.
    
    \vspace{0.3cm}
    Al inyectarlos en la ecuación dinámica, obtenemos un \textbf{perfil temporal de torque}[cite: 14]:
    \begin{equation*}
        \boxed{\tau(t) = M(q(t))\ddot{q}(t) + C(q(t), \dot{q}(t))\dot{q}(t) + g(q(t))} \text{[cite: 14]}
    \end{equation*}

    \vfill
    Esta formulación responde directamente a la pregunta de ingeniería: \textit{"¿Qué esfuerzo requerirán los motores a lo largo de toda la trayectoria planificada para el efector?"} integrando así la Cinemática (Unidad I) con la Dinámica (Unidad II).
\end{frame}

\begin{frame}{Integración: LSPB hacia Torque (2/2)[cite: 14]}
    Para nuestro ejemplo 1R y LSPB de clase ($q_i=0 \to q_f=2, t_f=3s, \ddot{q}_c=1$)[cite: 14]:
    
    \begin{center}
        \footnotesize
        \renewcommand{\arraystretch}{1.2}
        \begin{tabular}{l c c c c | c}
        \toprule
        \textbf{Fase} & $t$ (s) & $q$ (rad) & $\dot{q}$ (rad/s) & $\ddot{q}$ (rad/s$^2$) & $\tau$ \textit{(1R ilustrativo)} \\ \midrule
        Aceleración & $0.5$ & $0.125$ & $0.5$ & $1$ & $0.86$ N$\cdot$m \\
        Crucero & $1.5$ & $1.0$ & $1.0$ & $0$ & $4.13$ N$\cdot$m \\
        Desaceleración & $2.5$ & $1.875$ & $0.5$ & $-1$ & $4.43$ N$\cdot$m \\ \bottomrule
        \end{tabular}
    \end{center}
    
    \vfill
    \begin{alertblock}{Conclusión Fuerte}
        \small ¡Un tramo de aceleración nula ($\ddot{q}=0$) \textbf{no} significa torque cero! La gravedad o los términos dependientes de velocidad pueden dominar el esfuerzo requerido en ese instante[cite: 14].
    \end{alertblock}
\end{frame}

\begin{frame}{Singularidad Cinemática vs. Esfuerzo Dinámico[cite: 14]}
    \begin{columns}[T]
        \begin{column}{0.48\textwidth}
            \textbf{Singularidad (Cinemática)[cite: 14]:}
            
            Definida exclusivamente por la matriz geométrica:
            \begin{center}
                $\operatorname{rank}J(\mathbf{q}) < \operatorname{rank}_{\max}$ \text{[cite: 14]}
            \end{center}
            Implica la pérdida de grados de libertad operacionales instantáneos y posibles amplificaciones severas de velocidad articular en control por IK diferencial[cite: 14].
        \end{column}
        \begin{column}{0.48\textwidth}
            \textbf{Torque (Dinámica)[cite: 14]:}
            
            Definido por las inercias, velocidades y gravedad:
            \begin{center}
                $\tau(\mathbf{q}, \dot{\mathbf{q}}, \ddot{\mathbf{q}})$ \text{[cite: 14]}
            \end{center}
            \begin{tcolorbox}[colback=red!5, colframe=red!70!black, left=1mm, right=1mm]
                \small \textbf{Singularidad cinemática $\neq$ Torque alto}[cite: 14]. Una postura no singular puede exigir enorme torque gravitacional, y una postura singular se define por la movilidad, no por $\tau$[cite: 14].
            \end{tcolorbox}
        \end{column}
    \end{columns}
\end{frame}

\begin{frame}{Arquitectura Integrada del Hito 2[cite: 14]}
    El Hito 2 consolida toda la unidad en un solo pipeline de ingeniería (Python Notebook)[cite: 14]:
    
    \vspace{0.3cm}
    \begin{center}
    \begin{tikzpicture}[
        node distance=0.2cm and 0.5cm,
        box/.style={rectangle, draw=ucbblue, thick, fill=ucbblue!5, text width=2.4cm, align=center, rounded corners, font=\scriptsize},
        arrow/.style={-Stealth, thick, ucbblue}
    ]
        \node[box] (piep) {\textbf{Pieper}\\$T_d \to q^{(k)}$};
        \node[box, right=of piep] (fk) {\textbf{Verif. FK}\\$T_0^6(q) \approx T_d$};
        \node[box, right=of fk] (lspb) {\textbf{LSPB}\\$q(t), \dot{q}(t), \ddot{q}(t)$};
        
        \node[box, below=0.5cm of lspb] (dyn) {\textbf{Dinámica}\\$M, C, g \to \tau(t)$};
        \node[box, left=of dyn] (jac) {\textbf{Jacobiano}\\$J(q(t))$};
        \node[box, left=of jac] (sing) {\textbf{Singularidades}\\$\operatorname{rank}(J)$};
        
        \draw[arrow] (piep) -- (fk);
        \draw[arrow] (fk) -- (lspb);
        \draw[arrow] (lspb) -- (dyn);
        \draw[arrow] (lspb) -- (jac);
        \draw[arrow] (jac) -- (sing);
    \end{tikzpicture}
    \end{center}
    
    \vfill
    \begin{tcolorbox}[colback=white, colframe=ucbblue, title=Evidencia Final Esperada]
        \small $\boxed{ \text{Ecuaciones} + \text{Código} + \text{Gráficos} + \text{Validación} + \text{Interpretación} }$[cite: 14]
    \end{tcolorbox}
\end{frame}

\end{document}
